\section{Core Attack Mathematics}

To express the attack cleanly, define the \emph{intermediate value}

\begin{equation}
    I_i = D_K(C_i).
\end{equation}

Then CBC decryption becomes

\begin{equation}
    P_i = I_i \oplus C_{i-1}.
\end{equation}

The attacker does not know $I_i$, but the attacker can replace $C_{i-1}$ by a forged block $C'_{i-1}$. The receiver will then compute

\begin{equation}
    P'_i = I_i \oplus C'_{i-1}.
\end{equation}

The attack works by choosing $C'_{i-1}$ so that the last bytes of $P'_i$ become a valid PKCS\#7 padding. When the oracle returns ``valid,'' the attacker learns a relationship between the guessed bytes in $C'_{i-1}$ and the unknown bytes of $I_i$.

\subsection{Recovering the Last Byte}

Suppose the attacker wants the last byte of $P'_i$ to equal \texttt{01}. Let the attacker guess a value $g$ for the last byte of the forged previous block. The receiver computes

\[
P'_i[16] = I_i[16] \oplus g.
\]

If the oracle returns ``valid'' for a one-byte padding, then

\[
I_i[16] \oplus g = 01,
\]

which means

\begin{equation}
    I_i[16] = g \oplus 01.
\end{equation}

Once $I_i[16]$ is known, the last byte of the original plaintext follows immediately:

\begin{equation}
    P_i[16] = I_i[16] \oplus C_{i-1}[16].
\end{equation}

\subsection{Recovering Earlier Bytes}

After recovering the final byte, the attacker moves left. To recover byte 15, the attacker first forces the last byte of $P'_i$ to become \texttt{02} and then searches for a guess that also makes byte 15 equal \texttt{02}. More generally, when recovering byte position $j$, the attacker fixes all bytes to the right of $j$ so that they decrypt to the desired padding value

\[
16 - j + 1.
\]

The process repeats until the whole block is known. On average, each byte requires around 128 oracle queries because there are 256 possibilities for a byte and only one correct value is expected.

\subsection{Small Dry Run Example}

The abstract equations become much easier to follow with one fixed example. Assume AES has a 16-byte block size and the attacker submits \emph{two ciphertext blocks} to the oracle:

\[
C'_{i-1} \,\|\, C_i
\]

Here $C_i$ is the target block kept unchanged, while $C'_{i-1}$ is the forged previous block chosen by the attacker. The oracle decrypts $C_i$, XORs it with $C'_{i-1}$, and only tells the attacker whether the result ends with valid padding.

For explanation purposes, suppose the \emph{real} plaintext of the target block is

\[
P_i = \texttt{Meet me at dawn!}
\]

which is exactly 16 bytes, so it fills one block:

\begin{center}
\texttt{4D 65 65 74 20 6D 65 20 61 74 20 64 61 77 6E 21}
\end{center}

Assume the original previous ciphertext block ends with

\begin{center}
\texttt{... 5C A5}
\end{center}

so the last two plaintext bytes are

\begin{center}
\texttt{6E 21 = "n!"}
\end{center}

The attacker does not know these values during the real attack. They are shown here only to make the dry run understandable.

\begin{table}[H]
    \centering
    \caption{Setup for the dry-run example}
    \label{tab:dry-run-setup}
    \begin{tabular}{@{}p{0.28\textwidth}p{0.62\textwidth}@{}}
        \toprule
        Item & Value \\
        \midrule
        Block size & 16 bytes \\
        Ciphertext sent to oracle & forged block $C'_{i-1}$ followed by fixed block $C_i$ \\
        Text form of $P_i$ & \texttt{Meet me at dawn!} \\
        Last two bytes of $C_{i-1}$ & \texttt{5C A5} \\
        Last two bytes of $P_i$ & \texttt{6E 21} = \texttt{n!} \\
        \bottomrule
    \end{tabular}
\end{table}

\begin{table}[H]
    \centering
    \caption{CBC view of the same example}
    \label{tab:dry-run-diagram}
    \begin{tabular}{@{}p{0.22\textwidth}p{0.22\textwidth}p{0.16\textwidth}p{0.28\textwidth}@{}}
        \toprule
        Stage & Previous block & Target block & Result suffix \\
        \midrule
        Original message & \texttt{... 5C A5} & \texttt{$C_i$} & \texttt{6E 21} = \texttt{n!} \\
        Last-byte attack & \texttt{... ?? 85} & \texttt{$C_i$} & \texttt{?? 01} \\
        Two-byte attack & \texttt{... 30 86} & \texttt{$C_i$} & \texttt{02 02} \\
        \bottomrule
    \end{tabular}
\end{table}

\paragraph{What stays fixed and what changes?}
The target block $C_i$ never changes. The attacker only changes the previous block. Because CBC decryption uses

\[
P'_i = D_K(C_i) \oplus C'_{i-1},
\]

changing $C'_{i-1}$ changes the plaintext that appears after decryption. The oracle does not reveal the plaintext directly. It only says whether the last bytes form valid padding.

\paragraph{Step 1: Recover byte 16.}
To recover the last plaintext byte, the attacker wants the final decrypted byte to become \texttt{01}, because \texttt{01} is valid PKCS\#7 padding.

Suppose the attacker tries a forged last byte

\[
C'_{i-1}[16] = \texttt{85}.
\]

Assume the oracle now returns \texttt{valid}. That means the last byte of the decrypted block became \texttt{01}. In symbols,

\[
I_i[16] \oplus \texttt{85} = \texttt{01}.
\]

Therefore,

\[
I_i[16] = \texttt{85} \oplus \texttt{01} = \texttt{84}.
\]

Now the attacker recovers the original last plaintext byte using the \emph{original} previous block byte \texttt{A5}:

\[
P_i[16] = I_i[16] \oplus C_{i-1}[16]
        = \texttt{84} \oplus \texttt{A5}
        = \texttt{21}.
\]

Hex \texttt{21} is the character \texttt{!}. So the attacker has recovered the last byte of the real plaintext.

\paragraph{Step 2: Recover byte 15.}
Now the attacker wants the last \emph{two} bytes of the forged plaintext to become \texttt{02 02}.

The last byte is already solved, so the attacker first forces byte 16 to decrypt to \texttt{02}:

\[
C'_{i-1}[16] = I_i[16] \oplus \texttt{02}
             = \texttt{84} \oplus \texttt{02}
             = \texttt{86}.
\]

Next the attacker guesses byte 15 of the forged block. Suppose the oracle returns \texttt{valid} when

\[
C'_{i-1}[15] = \texttt{30}.
\]

Then byte 15 of the decrypted block must be \texttt{02}, so

\[
I_i[15] \oplus \texttt{30} = \texttt{02}.
\]

Thus

\[
I_i[15] = \texttt{30} \oplus \texttt{02} = \texttt{32}.
\]

Now the attacker uses the original byte \texttt{5C} from $C_{i-1}$:

\[
P_i[15] = I_i[15] \oplus C_{i-1}[15]
        = \texttt{32} \oplus \texttt{5C}
        = \texttt{6E}.
\]

Hex \texttt{6E} is the character \texttt{n}. So the attacker has now recovered the last two plaintext bytes:

\begin{center}
\texttt{6E 21 = "n!"}
\end{center}

This is the key idea of the attack. The attacker never decrypts $C_i$ directly and never learns the key. Instead, the attacker learns one byte at a time by forcing valid padding patterns and using XOR to recover the original plaintext.

\subsection{Complexity and Practical Meaning}

The attack does not recover the key, yet it still defeats confidentiality. For an $n$-block message, the attacker can recover the plaintext with a manageable number of oracle queries. The exact cost depends on implementation details and the amount of retry logic available to the attacker, but conceptually the attack is efficient enough to be practical when the target system accepts repeated modified ciphertexts.

This is the main mathematical lesson of the project: the attacker is not breaking AES. The attacker is exploiting the algebra of CBC decryption together with a validation side channel.
